\documentclass[a4paper,11pt]{report}

\usepackage[latin1]{inputenc}
%\usepackage[francais]{babel}
\usepackage[T1]{fontenc}
\usepackage{amsmath}
\usepackage{dsfont}
\usepackage{calc}
\usepackage{array}
\usepackage[a4paper]{geometry}
\usepackage{graphicx}
%\usepackage{bbold}
\usepackage{tensor}
\geometry{tmargin=1.5cm}
\geometry{lmargin=1.5cm}
\geometry{scale=0.8, nohead}

%\geometry{tmargin=1.5cm}
%\geometry{bmargin=2cm}
\geometry{scale=0.8, nohead}


\usepackage{titlesec}
% \titleformat{\chapter}[display]% OLD
%     {\normalfont\huge\bfseries}{\chaptertitlename\ \thechapter}{20pt}{\Huge}% OLD
% \titlespacing*{\chapter}{0pt}{50pt}{40pt}% OLD
\titleformat{\chapter}[display]% NEW
    {\fontfamily{ptm}\Large\bfseries\centering}{\chaptertitlename\ \thechapter}{5pt}{\Large}% NEW
\titlespacing*{\chapter}{0pt}{30pt}{20pt}% NEW
\renewcommand*\thesection{\arabic{section}}

\renewcommand{\vec}[1]{\mathbf{#1}}
\newcommand{\tens}[1]{\stackrel{\leftrightarrow}{#1}}
\usepackage{color}
\newcommand{\red}[1]{\textcolor{red}{#1}}

\newcommand{\vecp}[1]{\mathbf{#1}_{||}}
\newcommand{\macrofield}[5]{\big[\vec{#1}(\vec{#2_{||}},#3,#5 )\big]_{#4_{||}}}
\newcommand{\macrofieldlong}[5]{\big[{#1}^L(\vec{#2_{||}},#3,#5 )\big]_{#4_{||}}}

\newcommand{\bra}{\langle}
\newcommand{\ket}{\rangle}

%%%%%%%%%%%%%%%%%%%%%%%%%%%%%%%%%%%%%%%%%%%%%%%%%%%%%%%%%%%%%%%%%%%%%%%%%%%%%%%%%%%%%%%%%%%%%%%%%%%%%%%%%%%%%%%%%%%
%
%%%%%%%%%%%%%%%%%%%%%%%%%%%%%%%%%%%%%%%%%%%%%%%%%%%%%%%%%%%%%%%%%%%%%%%%%%%%%%%%%%%%%%%%%%%%%%%%%%%%%%%%%%%%%%%%%%%

\begin{document}

\chapter*{Viewing the world with two bands}

\section{Semiconductor Bloch Equations (SBE)}

In order to compare and understand the results presented by other group, we work in the framework of the semiconductor Bloch equations.

\subsection{TDDFT version of SBE}

In Ref.~\cite{PhysRevB.77.075204}, authors have showm that it is possible to formulate the TDDFT equations in the form of semiconductor Bloch equations.
Here we report the main equations, using the same notations than in the paper. 
We start with the time-dependent wave function, expressed in the basis of the ground-state wave functions
\begin{equation}
 \Psi^{v_i}_{\vec{k}}(\vec{r},t) = \sum_{j, \vec{q}}\Big[ c^{v_iv_j}_{\vec{kq}}(t)\phi^{v_j}_{\vec{q}}(\vec{r}) + c^{v_ic_j}_{\vec{kq}}(t)\phi^{c_j}_{\vec{q}}(\vec{r})\Big].\nonumber
\end{equation}
The one-body density matrix is then given by~\cite{PhysRevB.77.075204}
\begin{equation}
 \rho_{\vec{k};\vec{qp}}^{v_i;l_m\overline{l}_n}(t) = c^{v_il_m}_{\vec{kq}}(t)[c^{v_i\overline{l}_n}_{\vec{kp}}(t)]^*.
\end{equation}
This definition is usefull for defining occupation of the states, and is, in principle, exact as it is a reformulation of the time-dependent Kohn-Sham equations~\cite{PhysRevB.77.075204}.
The density matrix equation satifies the following equation of motion~\cite{PhysRevB.77.075204}
\begin{eqnarray}
 i\frac{\partial}{\partial t}\rho_{\vec{k};\vec{qp}}^{v_i;l_ml_n'}(t) &=& \left[ H(t),\rho\right]_{\vec{k};\vec{qp}}^{v_i;l_ml_n'} \nonumber\\
  &=& \sum{l_j'',\vec{q'}} \left[ H_{\vec{k};\vec{qq'}}^{l_ml_j''}(t) \rho_{\vec{k};\vec{q'p}}^{v_i;l_j''l_n'}(t) - \rho_{\vec{k};\vec{qq'}}^{v_i;l_ml_j''}(t)H_{\vec{k};\vec{q'p}}^{l_j''l_n'}(t)   \right],
\end{eqnarray}
where the Hamiltonian matrix elements $ H_{\vec{k};\vec{qq'}}^{l_ml_j''}(t)$ are
\begin{eqnarray}
  H_{\vec{k};\vec{qq'}}^{l_ml_j''}(t) &=& \int_{cell} d^3\vec{r} \Psi^{l_m*}_{\vec{q}}(\vec{r}) H(t)\Psi^{l_j''}_{\vec{q'}}(\vec{r}) \nonumber\\
   &=& E_{\vec{k}}^{l_m}\delta_{l_ml_n'}\delta_{\vec{kq}} + \vec{E}(t)d^{l_ml_n'}_{\vec{kq}} + V_{H,\vec{kq}}^{l_ml_n'}(t) +  + V_{xc,\vec{kq}}^{l_ml_n'}(t).
\end{eqnarray}
The dipole matrix elements are given by
\begin{equation}
 d^{l_ml_n'}_{\vec{kq}} = \int d^3\vec{r} \Psi^{l_m*}_{\vec{q}}(\vec{r})\vec{r}\Psi^{l_j''}_{\vec{q'}}(\vec{r}),
\end{equation}
and $V_{H,\vec{kq}}^{l_ml_n'}(t)$ and $V_{xc,\vec{kq}}^{l_ml_n'}(t)$ are defined as
\begin{equation}
  V_{H,\vec{kq}}^{l_ml_n'}(t) = \int d^3\vec{r} \Psi^{l_m*}_{\vec{q}}(\vec{r})\Big[V_{H}[n](\vec{r},t) - V_{H}[n](\vec{r},t_0) \Big] \Psi^{l_j''}_{\vec{q'}}(\vec{r}).
\end{equation}


\subsection{Two-band version of TDDFT SBE}

From the previous equations, it is possible to derive a two-band model.~\cite{PhysRevB.77.075204} We now assume that the physical properties are well describe using the  highest valence band and the lowest conduction band.
We also assume the dipole approximation for the external field, therefore we obtain that
\begin{equation}
 c^{vv}_\vec{kq}(t) = \delta_{\vec{kq}} c^{vv}_\vec{k}(t), \qquad c^{vc}_\vec{kq}(t) = \delta_{\vec{kq}} c^{vc}_\vec{k}(t),\nonumber
\end{equation}
which means that the optical transition have no momentum transfered.
Obviously, we have $\rho^{vv}_\vec{k}(t) + \rho^{cc}_\vec{k}(t) = 1$, and $\rho^{vc*}_\vec{k}(t)=\rho^{cv}_\vec{k}(t)$.

In Ref.~\cite{PhysRevB.77.075204}, authors obtained the TDDFT version of the SBE for a two-level system
\begin{subequations}
 \begin{align}
  \frac{\partial }{\partial t} \rho^{vv}_\vec{k}(t) =& -2\mathrm{Im}\{[E(t)d^{cv}_{\vec{k}} +  V_{H,\vec{k}}^{cv}(t) +  V_{xc,\vec{k}}^{cv}(t)]\rho^{vc}_\vec{k}(t) \},\\
  \frac{\partial }{\partial t} \rho^{vc}_\vec{k}(t) =& -i[E_\vec{k}^v-E_\vec{k}^c]\rho^{vc}_\vec{k}(t) - i[\rho^{cc}_\vec{k}(t)-\rho^{vv}_\vec{k}(t)][E(t)d^{vc}_{\vec{k}}+V_{H,\vec{k}}^{vc}(t) +  V_{xc,\vec{k}}^{vc}(t)] \nonumber \\
   &-i[V_{H,\vec{k}}^{vv}(t) +  V_{xc,\vec{k}}^{vv}(t) - V_{H,\vec{k}}^{cc}(t) -  V_{xc,\vec{k}}^{cc}(t)]\rho^{vc}_\vec{k}(t)
 \end{align}
 \label{eq:twolevel_TDDFT}
\end{subequations}


\subsection{Corkum's equations}

In Ref.~\cite{PhysRevLett.113.073901}, the two-level density matrix equations used are
\begin{subequations}
 \begin{align}
 \dot{\pi}(\vec{K},t) = -\frac{\pi(\vec{K},t)}{T_2} - i\Omega(\vec{K},t)w(\vec{K},t)e^{-iS(\vec{K},t)}, \\
 \dot{n}_m(\vec{K},t) = is_m\Omega^*(\vec{K},t)\pi(\vec{K},t)e^{iS(\vec{K},t)} + c.c.,
 \end{align}
  \label{eq:twolevel_Corkum}
\end{subequations}
where $n_m$ is the valence ($m=v$) and the conduction ($m=c$) band population, and $w=n_v-n_c$ is the population difference. Also, we have
$s_m=-1,1$, for $m=v,c$ respectively, and $\vec{K}=\vec{k}-\vec{A}(t)$. Further, $T_2$ accounts for the dephasing.
$S(\vec{K},t) = \int_{-\infty}^t E_{g}[\vec{K}+\vec{A}(t')]dt'$ is the classical action and $\Omega(\vec{K},t) = \vec{F}(t)\ldotp\vec{d}(\vec{k})$, is the Rabi frequency, with $\vec{F}(t) = \frac{d\vec{A}(t)}{dt}$ and $d(\vec{k})$ is the transition dipole moment 
\begin{equation}
 \vec{d}(\vec{k}) = i\int d^3\vec{x} u^*_{v,\vec{k}}(\vec{x})\nabla_{\vec{k}}u_{c,\vec{k}}(\vec{x}),
 \label{eq:dipole_Corkum}
\end{equation}
with the periodic part of the Bloch function $u_{c,\vec{k}}(\vec{x})$.

\subsection{Comparison between the two formalism}

In order to link the two formalism, we follow Ref.~\cite{PhysRevLett.113.073901} and we introduce the transformation
\begin{equation}
 c^{mm}_\vec{k}(t) = b^{m}_\vec{k}(t)\mathrm{exp}(-i\int_{-\infty}^t E_{\vec{k}}^m dt'),
\end{equation}
where $m$ can be $v$ or $c$.

We now define $n_m = |b^{m}|^2$ and $\pi = b^{c*} b^{v}$.
These two quantity are related to the previous one as 
\begin{eqnarray}
  n_m(\vec{k},t) = \rho^{mm}_\vec{k}(t), \nonumber\\
  \pi(\vec{k},t)e^{iS(\vec{k},t)} = \rho^{vc}_\vec{k}(t),  \nonumber
\end{eqnarray}

We directly obtain that $\left[\rho^{cc}_\vec{k}(t) - \rho^{vv}_\vec{k}(t)\right] = w(\vec{k},t)$.

Using this notations, we Eq.~\ref{eq:twolevel_TDDFT} reads as
\begin{subequations}
 \begin{align}
  \dot{n}_v(\vec{k},t) =& i[E(t)d^{cv}_{\vec{k}} +  V_{H,\vec{k}}^{cv}(t) +  V_{xc,\vec{k}}^{cv}(t)]\pi(\vec{k},t)e^{iS(\vec{k},t)} + c.c,\\
  \frac{\partial }{\partial t} [\pi(\vec{k},t)e^{iS(\vec{k},t)}] =& -i[E_\vec{k}^v-E_\vec{k}^c]\pi(\vec{k},t)e^{iS(\vec{k},t)} - i w(\vec{k},t)[E(t)d^{vc}_{\vec{k}}+V_{H,\vec{k}}^{vc}(t) +  V_{xc,\vec{k}}^{vc}(t)] \nonumber \\
   &-i[V_{H,\vec{k}}^{vv}(t) +  V_{xc,\vec{k}}^{vv}(t) - V_{H,\vec{k}}^{cc}(t) -  V_{xc,\vec{k}}^{cc}(t)]\pi(\vec{k},t)e^{iS(\vec{k},t)}
 \end{align}
 \label{eq:twolevel_TDDFT_2}
\end{subequations}

The second equation can be rewritten after simple algebra as
\begin{eqnarray}
   \dot{\pi}(\vec{k},t) &=& 2i[E_\vec{k}^c(t)-E_\vec{k}^v(t)]\pi(\vec{k},t)  - i w(\vec{k},t)[E(t)d^{vc}_{\vec{k}}+V_{H,\vec{k}}^{vc}(t) +  V_{xc,\vec{k}}^{vc}(t)]e^{-iS(\vec{k},t)} \nonumber \\
   &&-i[V_{H,\vec{k}}^{vv}(t) +  V_{xc,\vec{k}}^{vv}(t) - V_{H,\vec{k}}^{cc}(t) -  V_{xc,\vec{k}}^{cc}(t)]\pi(\vec{k},t)\nonumber
\end{eqnarray}

In order to link with the equation of  Ref.~\cite{PhysRevLett.113.073901}, we now introduce use the relation between the dipole matrix element and the dipole moment matrix element,
\begin{equation}
 d^{mn}_{\vec{k}} = i\partial_\vec{k} u_{m,\vec{k}}(t)\delta_{mn} + i\int d^3\vec{r} u^*_{m,\vec{k}}(\vec{x})\nabla_{\vec{k}}u_{n,\vec{k}}(\vec{x})\nonumber
\end{equation}
Therefore, we get that $d^{vc}_{\vec{k}}$ is directly replaced by $\vec{d}(\vec{k})$ in Corkum's notation.
Thus, we obtain that $E(t)d^{vc}_{\vec{k}}$ corresponds to $\Omega(\vec{k},t)$.

Putting everything together, we obtain that
\begin{subequations}
 \begin{align}
  \dot{\pi}(\vec{k},t) =& 2i[E_\vec{k}^c(t)-E_\vec{k}^v(t)]\pi(\vec{k},t)  - i w(\vec{k},t)[\Omega(\vec{k},t) +V_{H,\vec{k}}^{vc}(t) +  V_{xc,\vec{k}}^{vc}(t)]e^{-iS(\vec{k},t)} \nonumber \\
   &-i[V_{H,\vec{k}}^{vv}(t) +  V_{xc,\vec{k}}^{vv}(t) - V_{H,\vec{k}}^{cc}(t) -  V_{xc,\vec{k}}^{cc}(t)]\pi(\vec{k},t)\\
   \dot{n}_v(\vec{k},t) =& -i[\Omega(\vec{k},t) +  V_{H,\vec{k}}^{vc}(t) +  V_{xc,\vec{k}}^{vc}(t)]\pi(\vec{k},t)e^{iS(\vec{k},t)} + c.c,
 \end{align}
 \label{eq:twolevel_TDDFT_3}
\end{subequations}
These two equations correspond to the exact (in the TDDFT sense) version of the two-band model. \\

These two expressions can now be compared to Corkum's equations.
If we assume the independent-particle approximation, as done in Ref.~\cite{PhysRevLett.113.073901}, i.e. we put $V_{H}=V_{xc}=0$, we obtain from Eq.~\ref{eq:twolevel_TDDFT_3} 
\begin{subequations}
 \begin{align}
  \dot{\pi}(\vec{k},t) =& 2i[E_\vec{k}^c(t)-E_\vec{k}^v(t)]\pi(\vec{k},t)  - i \Omega(\vec{k},t)w(\vec{k},t)e^{-iS(\vec{k},t)} \\
   \dot{n}_v(\vec{k},t) =& -i\Omega(\vec{k},t)\pi(\vec{k},t)e^{iS(\vec{k},t)} + c.c,
 \end{align}
 \label{eq:twolevel_TDDFT_4}
\end{subequations}

\subsection{Remarks and comparison with other work}

\subsubsection{IPA versus frozen Hamiltonian}
It is not necessary to assume the IPA to recover the result of Ref.~\cite{PhysRevLett.113.073901}. If we assume a frozen Hamiltonian, i.e., $V_{Hxc}(t)=V_{Hxc}(t_0)$, we have $V_{H,\vec{k}}^{mn}(t) = V_{xc,\vec{k}}^{mn}(t)=0$ due to their definition.
\emph{Therefore, the TDDFT two-level equations with a frozen Hamiltonian approximation also lead the two-band equations of Ref.~\cite{PhysRevLett.113.073901}.} This gives more validity to this approach than previously thought. \\

\subsubsection{Dephasing and missing term}

By comparing with Eq.~\ref{eq:twolevel_Corkum}, we note that authors are missing a 
There is one term missing, even assuming the IPA, which is the term $ 2i[E_\vec{k}^c(t)-E_\vec{k}^v(t)]\pi(\vec{k},t)$ in their equations.
They added on the other time a dephasing term to account for effects wich are not present in the model, such as the phonons, defects, \textit{etc}.


% Interestingly, we obtain an exact expression for $T_2$, within the IPA, by comparing with with Eq.~\ref{eq:twolevel_Corkum}
% \begin{equation}
%  \frac{1}{T_2} = -2i[E_\vec{k}^c(t)-E_\vec{k}^v(t)].
% \end{equation}

\subsubsection{Comparison with the SBE of Ref.~\cite{hohenleutner2015real}}

In  Ref.~\cite{hohenleutner2015real}, the authors give the following SBEs (in their notations)

\begin{equation}
 \hbar\frac{\partial}{\partial t} f^{e_1}_k = -2\mathrm{Im}\left[ \sum_{e_2\neq e_1}d_{e_1e_2}(k)E(t)p_{k}^{e_2e_1} + \sum_{h_\lambda}  d_{e_1h_\lambda}(k)E(t)(p_k^{h_\lambda e_1})^*\right] + |e|E(t)\nabla_k f_k^{e_1} + \hbar \frac{\partial}{\partial t} f_{k}^{e_1}|_{\mathrm{relax}},
\end{equation}

\begin{eqnarray}
 \hbar\frac{\partial}{\partial t} p^{h_ie_j}_k = \left( \epsilon_k^{e_j} + \epsilon_k^{h_i} - i\frac{\hbar}{T_2} \right) p_{k}^{h_ie_j} - d_{e_jh_i}(k)E(t)\left(1-f_{k}^{e_j}-f_k^{h_i}\right)
 + i|e|E(t)\nabla_kp_k^{h_ie_j} \nonumber\\
  + E(t)\sum_{e_\lambda\neq e_j} [d_{e_\lambda h_i}(k)p_{k}^{e_\lambda e_j} - d_{e_je_\lambda}(k)p_{k}^{h_ie_\lambda}]
  + E(t)\sum_{h_\lambda\neq h_i} [d_{h_\lambda h_i}(k)p_{k}^{h_\lambda e_j} - d_{e_jh_\lambda}(k)p_{k}^{h_ih_\lambda}],
\end{eqnarray}
 with $f_k^\lambda$ the electron and hole occupations, and $p_k^{\lambda\lambda'}$ for $\lambda\neq\lambda'$ the microscopic polarization.~\cite{hohenleutner2015real} Here $d_{\lambda\lambda'}$ are the dipole-matrix elements.
 A phenomenological dephasing time $T_2$ is added to capture dominant effects of the high-order correlations. To account for the damping of currents in a realistic system, a phenomenological carrier relaxation $ \hbar \frac{\partial}{\partial t} f_{k}^{e_1}|_{\mathrm{relax}} = -\frac{1}{\tau}f_{k,A}^\lambda$ is included, which damps the anti-symmetric part $f_{k,A}^\lambda = \frac{1}{2}[f_k^{\lambda}-f_{-k}^{\lambda}]$ of the carrier distributions.~\cite{hohenleutner2015real}

 In order to compare with Corkum's model, we now restrict this model to two bands. As we have a valence and a conduction band, we note $f_k^{e_i}$ as  $f_c(\vec{k},t)$ and $f_k^{h_j}$ as $1-f_v(\vec{k},t)$. Similarly $\epsilon_k^{e_j}$ is denoted now $E_\vec{k}^c(t)$ and $\epsilon_k^{h_i}$ is denoted $-E_\vec{k}^v(t)$.
 The previous equations now reads (in atomic units)
 
 \begin{equation}
 \dot{f}_v(\vec{k},t) = -2\mathrm{Im}\left[ d_{vc}(k)E(t)p_k^{vc}(t)\right] + E(t)\nabla_k f_v(\vec{k},t) + \frac{\partial}{\partial t} f_{k}^{v}|_{\mathrm{relax}},
\end{equation}

\begin{eqnarray}
 \dot{p}^{vc}_k(t) = \left( E_\vec{k}^c(t) - E_\vec{k}^v(t) - i\frac{1}{T_2} \right) p_{k}^{vc}(t) - d_{cv}(k)E(t)\left(f_v(\vec{k},t)-f_c(\vec{k},t)\right)
 + iE(t)\nabla_k p_k^{vc}(t),
\end{eqnarray}

% In order to compare directly to Corkum's model, we use the following transformation
% \begin{eqnarray}
%  n_m(\vec{k},t) &=& f_c(\vec{k},t)\nonumber\\
%  \pi(\vec{k},t)e^{iS(\vec{k},t)} &=&  ip^{vc}_k(t)\nonumber
% \end{eqnarray}
% 
% We finally obtain that
% 
%  \begin{equation}
%  \dot{n}_v(\vec{k},t) = -2\mathrm{Im}\left[ \Omega(\vec{k},t) \pi(\vec{k},t)e^{-iS(\vec{k},t)} \right] + E(t)\nabla_k n_c(\vec{k},t) + \dot{n}_{k}^{v}|_{\mathrm{relax}},
% \end{equation}
% 
% \begin{eqnarray}
%  \dot{\pi}(\vec{k},t)e^{iS(\vec{k},t)} + \pi(\vec{k},t)e^{iS(\vec{k},t)} = i\left( 2E_\vec{k}^c(t) - 2E_\vec{k}^v(t) - i\frac{1}{T_2} \right) p_{k}^{vc}(t) + i \Omega(\vec{k},t)w(\vec{k},t)
%  - E(t)\nabla_k  (\pi(\vec{k},t)e^{iS(\vec{k},t)}),
% \end{eqnarray}
%  
%  where $w$ and $\Omega$ have been defined previously

\section{Two-band model in practice}

\subsection{Further approximation in Corkum's model}

\subsubsection{Dipole momentum matrix elements}

In Ref.~\cite{PhysRevLett.113.073901}, further approximation are involved in the model.\\
The dipole moment define by Eq.~\ref{eq:dipole_Corkum} is in fact choosen to be $k$-independent and to be given by
$\vec{d}(\vec{k}) = \vec{d_0}= (3.46,3.46,3.94)$.
Therefore, the Rabi frequency becomes $\vec{k}$-independent $\Omega(\vec{K},t) = \Omega(t) = \vec{F}(t)\ldotp\vec{d_0}$.

\subsubsection{Band structure}

The first approximation related to the band-structure is given by the approximation on the dispersion
\begin{equation}
 E_m(\vec{k}) = E_{m,x}(k_x) + E_{m,y}(k_y) + E_{m,z}(k_z).
 \label{eq:approx_disp}
\end{equation}

The bands are further approximated as cosine bands with parameters
\begin{subequations}
 \begin{align}
  E_{v,i}(k_i) = \sum_{j=0}^{\infty} \alpha^{j}_{v,i} \cos(j k_ia_i)\\
  E_{c,i}(k_i) = E_g + \sum_{j=0}^{\infty} \alpha^{j}_{c,i} \cos(j k_ia_i)
 \end{align}
 \label{eq:approx_bands}
\end{subequations}

For the $x$ direction, the maximum value for $j$ is 5, whereas for the other directions, $j$ is limited to 1.\cite{PhysRevB.91.064302} The values for $a_i$ are $(a_x,a_y,a_z) = (5.32,6.14,9.83)$.
We see directly that their is something wrong because Eq.~\ref{eq:approx_disp} together with Eqs.~\ref{eq:approx_bands} yield a band-gap of $3E_g$. There is a $\frac{1}{3}$ factor missing in Eq.~\ref{eq:approx_disp}.
In the following, we use instead
\begin{equation}
 E_m(\vec{k}) = \frac{E_{m,x}(k_x) + E_{m,y}(k_y) + E_{m,z}(k_z)}{3}.
 \label{eq:approx_disp2}
\end{equation}

The parameters for Eqs.~\ref{eq:approx_bands} are given in the Tab.~\ref{tab:bands_param}. The value of the band-gap is chose to be the experimental bandgap of $E_g=3.3eV$.\cite{PhysRevB.91.064302} This meens that the LDA band structure get a scissors correction in this approach.

\begin{table}
\centering
\begin{tabular}{lcc}
  \hline\hline
   & Valence band & Conduction band \\
  \hline
  $a_{0,x}$ &  -0.0928 &  0.0898 \\
  $a_{1,x}$ &   0.0705 & -0.0814 \\
  $a_{2,x}$ &   0.0200 & -0.0024 \\
  $a_{3,x}$ &  -0.0012 & -0.0048 \\
  $a_{4,x}$ &   0.0029 & -0.0003 \\
  $a_{5,x}$ &   0.0006 & -0.0009 \\
  $a_{0,y}$ &  -0.0307 &  0.1147 \\
  $a_{1,y}$ &   0.0307 & -0.1147 \\
  $a_{0,z}$ &  -0.0059 &  0.0435 \\
  $a_{1,z}$ &   0.0059 & -0.0435 \\
  \hline
\end{tabular}
\caption{Coefficients of the expansion in Eq.~\ref{eq:approx_bands}, taken from Ref.~\cite{PhysRevB.91.064302}}
\label{tab:bands_param}
\end{table}



\subsection{Putting everything together}
The laser field is defined  by 
\begin{equation}
 \vec{F}(t) = \vec{\hat{x}}F_0f(t),\nonumber
\end{equation}
where $f(t)$ is a gaussian enveloppe with a FWHM equal to 10 cycles and cosine carrier with frequency $\omega_0$. The $\vec{\hat{x}}$-axis is choosen to be parallel to the $\Gamma-M$ direction, whereas the $\vec{\hat{y}}$- and $\vec{\hat{z}}$ are respectively parallel to the $\Gamma-K$ and $\Gamma-A$ (optical axis) directions.
Therefore, we have that
\begin{equation}
 \Omega(\vec{K},t) = F_0f(t)d_{0,x} = F_0d_{0,x}\cos(\omega_0t)\exp(-\frac{(t-t_0)^2}{2\tau^2}),
\end{equation}
where $2\sqrt{2\ln{2}}\tau = 10\omega_0$ or $\tau = \frac{5\omega_0}{\sqrt{2\ln{2}}}$.


Therefore, from Eqs.~\ref{eq:twolevel_Corkum} we obtain the set of equations to be computed.
\begin{subequations}
 \begin{align}
 \dot{\pi}(\vec{K},t) = -\frac{\pi(\vec{K},t)}{T_2} - i F_0f(t)d_{0,x}[n_v(\vec{K},t)-n_c(\vec{K},t)]e^{-iS(\vec{K},t)}, \\
 \dot{n}_v(\vec{K},t) = 2 F_0f(t)d_{0,x}\mathrm{Im}\{\pi(\vec{K},t)e^{iS(\vec{K},t)}\}, \\
 \dot{n}_c(\vec{K},t) = -2 F_0f(t)d_{0,x}\mathrm{Im}\{\pi(\vec{K},t)e^{iS(\vec{K},t)}\},\\
 S(\vec{K},t) = \int_{-\infty}^t E_{g}[\vec{K}+\vec{A}(t')]dt'
 \end{align}
  \label{eq:twolevel_Corkum_real}
\end{subequations}

The polarisation if given by~\cite{PhysRevLett.113.073901}
\begin{equation}
\vec{p}(\vec{K},t) = \vec{d}(\vec{K}+A(t))\pi(\vec{K},t)e^{iS(\vec{K},t)}+c.c. \nonumber
\end{equation}
which turns out to be, accounting for previous approximations 
\begin{equation}
\vec{p}(\vec{K},t) = 2\vec{d_0}\mathrm{Re}\{\pi(\vec{K},t)e^{iS(\vec{K},t)}\}. \nonumber
\end{equation}

Finally, the interband and intraband contributions to the current are given by
\begin{subequations}
 \begin{align}
  \vec{j}_{ra}(t) =& \sum_{m=c,v} \int_{\overline{BZ}}\vec{v}_m[\vec{K}+\vec{A}(t)] n_m(\vec{K},t)d^3\vec{K},\\
  %
  \vec{j}_{er}(t) =& \frac{d}{dt}\int_{\overline{BZ}} \vec{p}(\vec{K},t)d^3\vec{K},
 \end{align}
\end{subequations}
with $\vec{\vec{v}_m}(\vec{k}) = \nabla_{\vec{k}}E_m(\vec{k})$ and $\overline{BZ}=BZ-\vec{A}(t)$.

From the previous definitions, we obtain that the interband term reads

\begin{eqnarray}
  j_{er,i}(t) &=& 2d_{0,i}\mathrm{Re}\int_{\overline{BZ}}\Bigg(\dot{\pi}(\vec{K},t)+i\pi(\vec{K},t)\dot{S}(\vec{K},t)\Bigg) e^{iS(\vec{K},t)}d^3\vec{K}\} \nonumber\\
   &=& 2d_{0,i}\mathrm{Re}\{\int_{\overline{BZ}}\Bigg(\dot{\pi}(\vec{K},t)+i\pi(\vec{K},t)E_g[\vec{K}+\vec{A}(t)]\Bigg) e^{iS(\vec{K},t)}d^3\vec{K}\} \nonumber\\
%
%
   &=& -2iF_0f(t)d_{0,i} d_{0,x}\mathrm{Re}\{\int_{\overline{BZ}}[n_v(\vec{K},t)-n_c(\vec{K},t)]d^3\vec{K}\} \nonumber\\
   &&+2d_{0,i}\mathrm{Re}\{\int_{\overline{BZ}} \Big(iE_g[\vec{K}+\vec{A}(t)]-\frac{1}{T_2}\Big)\pi(\vec{K},t) e^{iS(\vec{K},t)}d^3\vec{K}\} \nonumber\\
\end{eqnarray}

whereas the intraband part reads, from Eq.~\ref{eq:approx_disp2} and Eq.~\ref{eq:approx_bands},
\begin{eqnarray}
 j_{ra,i}(t) &=& \sum_{m=v,c}\sum_{l=1}^{\infty} l \alpha^{l}_{m,i} \int_{BZ}
   n_m(\vec{k}-\vec{A}(t),t)k_i \sin(lk_ia_i)  d^3\vec{k}.
\end{eqnarray}

\subsection{Keldysh approximation}

In Ref.~\cite{PhysRevLett.113.073901} authors used the Keldysh approximation, \textit{i.e}, $w\approx 1$ to decouples the equations.
Doing that, we obtain that
\begin{subequations}
 \begin{align}
 \dot{\pi}(\vec{K},t) = -\frac{\pi(\vec{K},t)}{T_2} - i F_0f(t)d_{0,x}e^{-iS(\vec{K},t)}, \\
 \dot{n}_m(\vec{K},t) = 2s_m F_0f(t)d_{0,x}\mathrm{Im}\{\pi(\vec{K},t)e^{iS(\vec{K},t)}\},
 \end{align}
\end{subequations}

Let us first write the first equation as
\begin{equation}
 \dot{y}(t) = by(t) + c(t),
\end{equation}
with
\begin{eqnarray}
 y(t) &=& \pi(\vec{K},t)\nonumber\\
 b &=& -\frac{1}{T_2} \nonumber\\
 c(t) &=& - i F_0f(t)d_{0,x}e^{-iS(\vec{K},t)}\nonumber
\end{eqnarray}
We have
\begin{eqnarray}
 y(t) = \lambda(t)e^{G(t)},\\
 \dot{\lambda}(t) = c(t)e^{-G(t)}\\
 \dot{G}(t) = -b = \frac{1}{T_2} \rightarrow G(t) = \frac{t}{T_2}+C^{ste}\\
\end{eqnarray}
so
\begin{equation}
  \dot{\lambda}(t)  =  i\vec{d_{0}}\vec{F}(t)e^{-(iS(\vec{K},t)+\frac{t}{T_2})}
\end{equation}

\begin{eqnarray}
  \lambda(t_b)-\lambda(t_a) = i\vec{d_{0}}\int_{t_a}^{t_b}\vec{F}(t)e^{-(iS(\vec{K},t)+\frac{t}{T_2})}dt% \\
 % =  i\vec{d_{0}} \Bigg[ \vec{A}(t)e^{-(iS(\vec{K},t)+\frac{t}{T_2})}\Bigg]_{t_a}^{t_b} - i\vec{d_{0}} \int_{t_a}^{t_b}\vec{A}(t)[-iE_g(\vec{K},t)+\frac{1}{T_2}]e^{-(iS(\vec{K},t)+\frac{t}{T_2})}dt
\end{eqnarray}

Using the condition $y(t=-\infty)=0$ and $\vec{A}(t=-\infty)=0$, we obtain that
% \begin{eqnarray}
%   \lambda(t) = i\int_{-\infty}^{t}\vec{d_{0}}\ldotp\vec{F}(t')e^{-(iS(\vec{K},t')+\frac{t'}{T_2})}dt'
% \end{eqnarray}
% % \begin{eqnarray}
% %   \lambda(t) =  i\vec{d_{0}} \vec{A}(t)e^{-(iS(\vec{K},t)+\frac{t}{T_2})} - i\vec{d_{0}} \int_{-\infty}^{t}\vec{A}(t')\left[-iE_g(\vec{K},t')+\frac{1}{T_2}\right]e^{-(iS(\vec{K},t')+\frac{t'}{T_2})}dt'
% % \end{eqnarray}
% 
% which leads to the solution
\begin{subequations}
 \begin{align}
 \pi(\vec{K},t) =& \int_{-\infty}^{t}\vec{d_{0}}\ldotp\vec{F}(t')e^{-iS(\vec{K},t')+\frac{t-t'}{T_2})}dt', \\
 n_m(\vec{K},t) =& 2s_m F_0f(t)d_{0,x}\mathrm{Im}\{\int_{-\infty}^{t}\vec{d_{0}}\ldotp\vec{F}(t')e^{iS(\vec{K},t)-iS(\vec{K},t')+\frac{t-t'}{T_2})}dt'\} . 
 \end{align}
\end{subequations}
% \begin{subequations}
%  \begin{align}
%  \pi(\vec{K},t) =& i\vec{d_{0}} \vec{A}(t)e^{-iS(\vec{K},t)} - \vec{d_{0}} \int_{-\infty}^{t}\vec{A}(t')\left[E_g(\vec{K},t')+\frac{i}{T_2}\right]e^{-(iS(\vec{K},t')+\frac{(t'-t)}{T_2})}dt', \\
%  n_m(\vec{K},t) =& 2s_m F_0f(t)d_{0,x}\mathrm{Im}\Big\{i\vec{d_{0}} \vec{A}(t) - \vec{d_{0}} \int_{-\infty}^{t}\vec{A}(t')\left[E_g(\vec{K},t')+\frac{i}{T_2}\right]e^{iS(\vec{K},t)-iS(\vec{K},t')+\frac{(t-t')}{T_2})}dt'\Big\}. \nonumber\\
%  =&  2s_m F_0f(t)d_{0,x}\Bigg(\vec{d_{0}} \vec{A}(t)  
%  - \mathrm{Im}\Big\{\vec{d_{0}} \int_{-\infty}^{t}\vec{A}(t')\left[E_g(\vec{K},t')+\frac{i}{T_2}\right]e^{iS(\vec{K},t)-iS(\vec{K},t')+\frac{(t-t')}{T_2})}dt'\Big\}\Bigg). 
%  \end{align}
% \end{subequations}
% 
% Using the fact that for $t'<t$, we have
% \begin{eqnarray}
%  S(\vec{K},t)-S(\vec{K},t') = \int_{-\infty}^t E_{g}[\vec{K}+\vec{A}(t'')]dt'' - \int_{-\infty}^{t'} E_{g}[\vec{K}+\vec{A}(t'')]dt'' \nonumber\\
%   = \int_{t'}^{t} E_{g}[\vec{K}+\vec{A}(t'')]dt''
% \end{eqnarray}
% This quantity is now called $S(\vec{K},t',t)$, following Ref.~\cite{PhysRevLett.113.073901}.
% 
% Therefore, 
% \begin{equation}
%  n_m(\vec{K},t)  =  2s_m F_0f(t)d_{0,x}\Bigg(\vec{d_{0}} \vec{A}(t)  
%  - \mathrm{Im}\Big\{\vec{d_{0}} \int_{-\infty}^{t}\vec{A}(t')\left[E_g(\vec{K},t')+\frac{i}{T_2}\right]e^{iS(\vec{k},t',t)+\frac{(t-t')}{T_2})}dt'\Big\}\Bigg). 
% \end{equation}

Inserting these equations into the expressions of the interband and intraband current yields
\begin{subequations}
 \begin{align}
  \vec{j}_{ra}(t) =& \sum_{m=c,v} 2s_m \int_{BZ}d^3\vec{k}\vec{v}_m(\vec{k}) \nonumber\\
   &\times \vec{d_{0}}\ldotp\vec{F}(t)\mathrm{Im}\{\int_{-\infty}^{t}\vec{d_{0}}\ldotp\vec{F}(t')e^{iS(\vec{k},t',t)+\frac{t-t'}{T_2})}dt'\},\\
  %
  \vec{j}_{er}(t) =& \frac{d}{dt}\int_{BZ}d^3\vec{k}  \vec{d_0}\int_{-\infty}^{t}dt'\vec{d_{0}}\ldotp\vec{F}(t')2\mathrm{Re}\{e^{iS(\vec{K},t',t)+\frac{t-t'}{T_2})}\},
 \end{align}
\end{subequations}
where $S(\vec{K},t',t)$ is defined following Ref.~\cite{PhysRevLett.113.073901}.
Fourier tranform of these equations gives the equations of Ref.~\cite{PhysRevLett.113.073901}.

% 
% \begin{eqnarray}
%  j_{ra,i}(t) &=& -\sum_{m=c,v} 2s_m F_0f(t)d_{0,x}\int_{BZ}\vec{v}_m(\vec{k}) 
%  \mathrm{Im}\Big\{\int_{-\infty}^{t}\vec{d_{0}} \vec{A}(t')\left[E_g(\vec{K},t')+\frac{i}{T_2}\right]e^{iS(\vec{k}-\vec{A}(t),t',t)+\frac{(t-t')}{T_2})}dt'\Big\}\Bigg)d^3\vec{k}.
% \end{eqnarray}
% 
% 
% \begin{eqnarray}
%   j_{er,i}(t) &=&-2d_{0,i}\vec{d_{0}} \vec{A}(t)\mathrm{Re}\{\int_{\overline{BZ}} \Big(E_g[\vec{K}+\vec{A}(t)]+\frac{i}{T_2}\Big)  d^3\vec{K}\} \nonumber\\
%   &&
%   - 2d_{0,i}\mathrm{Re}\{\int_{\overline{BZ}} \Big(iE_g[\vec{K}+\vec{A}(t)]-\frac{1}{T_2}\Big)\Bigg[\vec{d_{0}} \int_{t_0}^{t}\vec{A}(t')\left[E_g(\vec{K},t')+\frac{i}{T_2}\right]e^{-(iS(\vec{K},t')+\frac{(t'-t)}{T_2})}dt' e^{iS(\vec{K},t)}d^3\vec{K}\} 
% \end{eqnarray}
% 
% \begin{eqnarray}
%   j_{er,i}(t) &=&-2d_{0,i}\vec{d_{0}} \vec{A}(t)\int_{BZ}E_g[\vec{k}]d^3\vec{k} \nonumber\\
%   &&
%   - 2d_{0,i}\mathrm{Re}\{\int_{\overline{BZ}} \Big(iE_g[\vec{K}+\vec{A}(t)]-\frac{1}{T_2}\Big)\Bigg[\vec{d_{0}} \int_{t_0}^{t}\vec{A}(t')\left[E_g(\vec{K},t')+\frac{i}{T_2}\right] e^{iS(\vec{K},t)-iS(\vec{K},t')+\frac{(t-t')}{T_2})}dt'd^3\vec{K}\} 
% \end{eqnarray}


\bibliographystyle{unsrt}
\bibliography{./biblio}
\end{document}